\documentclass[10pt,a4paper]{amsart}
\usepackage{amsmath,amssymb,amsthm,graphicx}
\usepackage{verbatim}
\newenvironment{reference}{\comment}{\endcomment}
%\newenvironment{reference}{}{}
\newenvironment{slogan}{\comment}{\endcomment}
\newenvironment{history}{\comment}{\endcomment}

% For commutative diagrams we use Xy-pic
\usepackage[all]{xy}

% We use 2cell for 2-commutative diagrams.
\xyoption{2cell}
% Theorem environments.
%
\theoremstyle{plain}
\newtheorem{theorem}[subsection]{Theorem}
\newtheorem{proposition}[subsection]{Proposition}
\newtheorem{lemma}[subsection]{Lemma}

\theoremstyle{definition}
\newtheorem{definition}[subsection]{Definition}
\newtheorem{example}[subsection]{Example}
\newtheorem{exercise}[subsection]{Exercise}
\newtheorem{situation}[subsection]{Situation}

\theoremstyle{remark}
\newtheorem{remark}[subsection]{Remark}
\newtheorem{remarks}[subsection]{Remarks}

\numberwithin{equation}{subsection}

% Macros
%
\def\lim{\mathop{\mathrm{lim}}\nolimits}
\def\colim{\mathop{\mathrm{colim}}\nolimits}
\def\Spec{\mathop{\mathrm{Spec}}}
\def\Hom{\mathop{\mathrm{Hom}}\nolimits}
\def\Ext{\mathop{\mathrm{Ext}}\nolimits}
\def\SheafHom{\mathop{\mathcal{H}\!\mathit{om}}\nolimits}
\def\SheafExt{\mathop{\mathcal{E}\!\mathit{xt}}\nolimits}
\def\Sch{\mathit{Sch}}
\def\Mor{\mathop{\mathrm{Mor}}\nolimits}
\def\Ob{\mathop{\mathrm{Ob}}\nolimits}
\def\Sh{\mathop{\mathit{Sh}}\nolimits}
\def\NL{\mathop{N\!L}\nolimits}
\def\CH{\mathop{\mathrm{CH}}\nolimits}
\def\proetale{{pro\text{-}\acute{e}tale}}
\def\etale{{\acute{e}tale}}
\def\QCoh{\mathit{QCoh}}
\def\Ker{\mathop{\mathrm{Ker}}}
\def\Im{\mathop{\mathrm{Im}}}
\def\Coker{\mathop{\mathrm{Coker}}}
\def\Coim{\mathop{\mathrm{Coim}}}

% Boxtimes
%
\DeclareMathSymbol{\boxtimes}{\mathbin}{AMSa}{"02}

%
% Macros for moduli stacks/spaces
%
\def\QCohstack{\mathcal{QC}\!\mathit{oh}}
\def\Cohstack{\mathcal{C}\!\mathit{oh}}
\def\Spacesstack{\mathcal{S}\!\mathit{paces}}
\def\Quotfunctor{\mathrm{Quot}}
\def\Hilbfunctor{\mathrm{Hilb}}
\def\Curvesstack{\mathcal{C}\!\mathit{urves}}
\def\Polarizedstack{\mathcal{P}\!\mathit{olarized}}
\def\Complexesstack{\mathcal{C}\!\mathit{omplexes}}
% \Pic is the operator that assigns to X its picard group, usage \Pic(X)
% \Picardstack_{X/B} denotes the Picard stack of X over B
% \Picardfunctor_{X/B} denotes the Picard functor of X over B
\def\Pic{\mathop{\mathrm{Pic}}\nolimits}
\def\Picardstack{\mathcal{P}\!\mathit{ic}}
\def\Picardfunctor{\mathrm{Pic}}
\def\Deformationcategory{\mathcal{D}\!\mathit{ef}}
\usepackage{cleveref}
\crefname{originallemma}{original lemma}{original lemmas}
\Crefname{originallemma}{Original lemma}{Original lemmas}
\makeatletter
\newlabel{descent-lemma-descending-property-finite}{{02LA}{source}{}{sourcepointer-1}{}}
\newlabel{descent-lemma-descending-property-finite@cref}{{[originallemma][0][]02LA}{[1][0][]source}{}{}{}}
\makeatother
\crefname{originaldefinition}{original definition}{original definitions}
\Crefname{originaldefinition}{Original definition}{Original definitions}
\makeatletter
\newlabel{groupoids-definition-invariant-open}{{03BC}{source}{}{sourcepointer-2}{}}
\newlabel{groupoids-definition-invariant-open@cref}{{[originaldefinition][0][]03BC}{[1][0][]source}{}{}{}}
\makeatother
\crefname{originalreference}{original reference}{original references}
\Crefname{originalreference}{Original reference}{Original references}
\makeatletter
\newlabel{lemma-quasi-separated-finite-etale-cover-dense-open-scheme}{{0BAS}{source}{}{sourcepointer-3}{}}
\newlabel{lemma-quasi-separated-finite-etale-cover-dense-open-scheme@cref}{{[originalreference][0][]0BAS}{[1][0][]source}{}{}{}}
\makeatother
\crefname{originalreference}{original reference}{original references}
\Crefname{originalreference}{Original reference}{Original references}
\makeatletter
\newlabel{lemma-quasi-separated-quasi-compact-pieces}{{03W7}{source}{}{sourcepointer-4}{}}
\newlabel{lemma-quasi-separated-quasi-compact-pieces@cref}{{[originalreference][0][]03W7}{[1][0][]source}{}{}{}}
\makeatother
\crefname{originalreference}{original reference}{original references}
\Crefname{originalreference}{Original reference}{Original references}
\makeatletter
\newlabel{lemma-separated-cover}{{03FZ}{source}{}{sourcepointer-5}{}}
\newlabel{lemma-separated-cover@cref}{{[originalreference][0][]03FZ}{[1][0][]source}{}{}{}}
\makeatother
\crefname{originalreference}{original reference}{original references}
\Crefname{originalreference}{Original reference}{Original references}
\makeatletter
\newlabel{lemma-subscheme}{{03JH}{source}{}{sourcepointer-6}{}}
\newlabel{lemma-subscheme@cref}{{[originalreference][0][]03JH}{[1][0][]source}{}{}{}}
\makeatother
\crefname{originallemma}{original lemma}{original lemmas}
\Crefname{originallemma}{Original lemma}{Original lemmas}
\makeatletter
\newlabel{more-groupoids-lemma-property-invariant}{{03JC}{source}{}{sourcepointer-7}{}}
\newlabel{more-groupoids-lemma-property-invariant@cref}{{[originallemma][0][]03JC}{[1][0][]source}{}{}{}}
\makeatother
\crefname{originallemma}{original lemma}{original lemmas}
\Crefname{originallemma}{Original lemma}{Original lemmas}
\makeatletter
\newlabel{morphisms-lemma-etale-locally-quasi-finite}{{03WS}{source}{}{sourcepointer-8}{}}
\newlabel{morphisms-lemma-etale-locally-quasi-finite@cref}{{[originallemma][0][]03WS}{[1][0][]source}{}{}{}}
\makeatother
\crefname{originallemma}{original lemma}{original lemmas}
\Crefname{originallemma}{Original lemma}{Original lemmas}
\makeatletter
\newlabel{morphisms-lemma-generically-finite}{{02NW}{source}{}{sourcepointer-9}{}}
\newlabel{morphisms-lemma-generically-finite@cref}{{[originallemma][0][]02NW}{[1][0][]source}{}{}{}}
\makeatother
\crefname{originallemma}{original lemma}{original lemmas}
\Crefname{originallemma}{Original lemma}{Original lemmas}
\makeatletter
\newlabel{morphisms-lemma-quasi-finite}{{02NH}{source}{}{sourcepointer-10}{}}
\newlabel{morphisms-lemma-quasi-finite@cref}{{[originallemma][0][]02NH}{[1][0][]source}{}{}{}}
\makeatother
\crefname{originallemma}{original lemma}{original lemmas}
\Crefname{originallemma}{Original lemma}{Original lemmas}
\makeatletter
\newlabel{morphisms-lemma-quasi-finite-locally-quasi-compact}{{01TJ}{source}{}{sourcepointer-11}{}}
\newlabel{morphisms-lemma-quasi-finite-locally-quasi-compact@cref}{{[originallemma][0][]01TJ}{[1][0][]source}{}{}{}}
\makeatother
\crefname{originallemma}{original lemma}{original lemmas}
\Crefname{originallemma}{Original lemma}{Original lemmas}
\makeatletter
\newlabel{spaces-lemma-finding-opens}{{02WU}{source}{}{sourcepointer-12}{}}
\newlabel{spaces-lemma-finding-opens@cref}{{[originallemma][0][]02WU}{[1][0][]source}{}{}{}}
\makeatother
\title{A Zariski locally quasi-separated algebraic space contains a dense open subspace which is a scheme}
\author{The Stacks Project Authors}
\date{stacks-project:spaces-properties.tex:local-snapshot-20261012; GNU FDL 1.2 or later}
\begin{document}
\maketitle
\noindent\textit{Editorial note.} All delivered mathematical text is the author's own, copied from the pinned Stacks Project chapter file spaces-properties.tex; no statement or proof step was rewritten, reordered or removed.\par\medskip
\noindent\textit{Editorial note.} The counted proof is the author's complete argument: the reduction to a single affine chart with the properties of the two projections, the construction of the invariant open over which the source projection is finite by generic finiteness and descent, and the passage to the quotient of that open reducing the statement to the finite etale case.\par\medskip
\noindent\textit{Editorial note.} The label proposition-locally-quasi-separated-open-dense-scheme names the counted statement, and the proof environment immediately following it in the source is delivered verbatim as its proof.\par\medskip
\noindent\textit{Editorial note.} The counted statement is carried by the author's own proposition environment in the chapter source.\par\medskip
\noindent\textit{Editorial note.} Pre-open gate: python operations/agent-authoring/tools/precard\_gate.py --source spaces-properties.tex --label proposition-locally-quasi-separated-open-dense-scheme --start 1863 --end 1906 --local-id 28048 returned RESULT=OK (L28048 free, occupancy FREE, public interval clear). Candidate taken from sources/stacks-spaces-properties/candidates-halfA.json.\par\medskip
\noindent\textit{Editorial note.} The statement and proof are copied from the pinned chapter file; the cited inputs are kept as external pointers with their canonical Stacks tags.\par\medskip
\noindent\textit{Allowed standard prerequisites (editorial).}\begin{itemize}
\item Local nature of the statement in the space. Same chapter, the lemma that the property of containing a dense open subscheme can be checked Zariski locally. Supplies the reduction to a single piece of a Zariski cover.
\item Characters of the charts given by the local quasi-separatedness. Same chapter, the lemma characterizing Zariski local quasi-separatedness through quasi-compact surjective etale morphisms from affine schemes. Supplies the affine chart U -> X, which is separated because the cover is separated.
\item Generic finiteness, descent of finiteness and invariance of the locus of finiteness. Morphisms, the lemma producing an open neighbourhood of a generic point over which a quasi-finite morphism becomes finite, Descent, the lemma that finiteness descends along fpqc morphisms, and More on Groupoids, the lemma that the largest open locus where the source map is finite is invariant. Supply the invariant open W containing all generic points of irreducible components of U, over which the relation is finite.
\item Restriction of the relation to an invariant open and the finite etale case. Groupoids, the description of the restriction of a groupoid to an invariant open, and Same chapter, the lemma treating the case of a finite etale cover. Supply the open subspace X' = W/R\_W and the final reduction to the finite etale case.
\end{itemize}\medskip
\ifdefined\hypertarget\hypertarget{sourcepointer-1}{}\fi
\noindent\textit{Source reference.} Original lemma 02LA: sources/stacks-chow/tags.txt, source lines 3401--3401. Prerequisite pointer invoked at proof line(s) 1889; it names a standard fact of a sibling chapter that the counted argument uses, and every step of the counted proof built on it is retained verbatim in the delivered body.\par\medskip
\ifdefined\hypertarget\hypertarget{sourcepointer-2}{}\fi
\noindent\textit{Source reference.} Original definition 03BC: sources/stacks-chow/tags.txt, source lines 4343--4343. Prerequisite pointer invoked at proof line(s) 1897; it names a standard fact of a sibling chapter that the counted argument uses, and every step of the counted proof built on it is retained verbatim in the delivered body.\par\medskip
\ifdefined\hypertarget\hypertarget{sourcepointer-3}{}\fi
\noindent\textit{Source reference.} Original reference 0BAS: sources/stacks-spaces-properties/pinned-original-source.tex, source lines 1817--1817. Prerequisite pointer invoked at proof line(s) 1905; it names a standard fact of the same chapter that the counted argument uses, and every step of the counted proof built on it is retained verbatim in the delivered body.\par\medskip
\ifdefined\hypertarget\hypertarget{sourcepointer-4}{}\fi
\noindent\textit{Source reference.} Original reference 03W7: sources/stacks-spaces-properties/pinned-original-source.tex, source lines 1065--1065. Prerequisite pointer invoked at proof line(s) 1874; it names a standard fact of the same chapter that the counted argument uses, and every step of the counted proof built on it is retained verbatim in the delivered body.\par\medskip
\ifdefined\hypertarget\hypertarget{sourcepointer-5}{}\fi
\noindent\textit{Source reference.} Original reference 03FZ: sources/stacks-spaces-properties/pinned-original-source.tex, source lines 992--992. Prerequisite pointer invoked at proof line(s) 1877; it names a standard fact of the same chapter that the counted argument uses, and every step of the counted proof built on it is retained verbatim in the delivered body.\par\medskip
\ifdefined\hypertarget\hypertarget{sourcepointer-6}{}\fi
\noindent\textit{Source reference.} Original reference 03JH: sources/stacks-spaces-properties/pinned-original-source.tex, source lines 1782--1782. Prerequisite pointer invoked at proof line(s) 1873; it names a standard fact of the same chapter that the counted argument uses, and every step of the counted proof built on it is retained verbatim in the delivered body.\par\medskip
\ifdefined\hypertarget\hypertarget{sourcepointer-7}{}\fi
\noindent\textit{Source reference.} Original lemma 03JC: sources/stacks-chow/tags.txt, source lines 4634--4634. Prerequisite pointer invoked at proof line(s) 1892; it names a standard fact of a sibling chapter that the counted argument uses, and every step of the counted proof built on it is retained verbatim in the delivered body.\par\medskip
\ifdefined\hypertarget\hypertarget{sourcepointer-8}{}\fi
\noindent\textit{Source reference.} Original lemma 03WS: sources/stacks-chow/tags.txt, source lines 5123--5123. Prerequisite pointer invoked at proof line(s) 1882; it names a standard fact of a sibling chapter that the counted argument uses, and every step of the counted proof built on it is retained verbatim in the delivered body.\par\medskip
\ifdefined\hypertarget\hypertarget{sourcepointer-9}{}\fi
\noindent\textit{Source reference.} Original lemma 02NW: sources/stacks-chow/tags.txt, source lines 3495--3495. Prerequisite pointer invoked at proof line(s) 1885; it names a standard fact of a sibling chapter that the counted argument uses, and every step of the counted proof built on it is retained verbatim in the delivered body.\par\medskip
\ifdefined\hypertarget\hypertarget{sourcepointer-10}{}\fi
\noindent\textit{Source reference.} Original lemma 02NH: sources/stacks-chow/tags.txt, source lines 3480--3480. Prerequisite pointer invoked at proof line(s) 1884; it names a standard fact of a sibling chapter that the counted argument uses, and every step of the counted proof built on it is retained verbatim in the delivered body.\par\medskip
\ifdefined\hypertarget\hypertarget{sourcepointer-11}{}\fi
\noindent\textit{Source reference.} Original lemma 01TJ: sources/stacks-chow/tags.txt, source lines 2376--2376. Prerequisite pointer invoked at proof line(s) 1883; it names a standard fact of a sibling chapter that the counted argument uses, and every step of the counted proof built on it is retained verbatim in the delivered body.\par\medskip
\ifdefined\hypertarget\hypertarget{sourcepointer-12}{}\fi
\noindent\textit{Source reference.} Original lemma 02WU: sources/stacks-chow/tags.txt, source lines 3817--3817. Prerequisite pointer invoked at proof line(s) 1901; it names a standard fact of a sibling chapter that the counted argument uses, and every step of the counted proof built on it is retained verbatim in the delivered body.\par\medskip
\begin{proposition}
\label{proposition-locally-quasi-separated-open-dense-scheme}
Let $S$ be a scheme. Let $X$ be an algebraic space over $S$. If $X$ is
Zariski locally quasi-separated (for example if $X$ is quasi-separated), then
there exists a dense open subspace $X'$ of $X$ which is a scheme.
More precisely, every point $x \in |X|$ of codimension $0$ on $X$
is contained in $X'$.
\end{proposition}

\begin{proof}
The question is local on $X$ by Lemma \ref{lemma-subscheme}.
Thus by Lemma \ref{lemma-quasi-separated-quasi-compact-pieces}
we may assume that there exists an affine scheme $U$ and a
surjective, quasi-compact, \'etale morphism $U \to X$.
Moreover $U \to X$ is separated (Lemma \ref{lemma-separated-cover}).
Set $R = U \times_X U$ and denote $s, t : R \to U$ the projections
as usual. Then $s, t$ are surjective, quasi-compact, separated, and
\'etale. Hence $s, t$ are also quasi-finite and have finite fibres
(Morphisms,
Lemmas \ref{morphisms-lemma-etale-locally-quasi-finite},
\ref{morphisms-lemma-quasi-finite-locally-quasi-compact}, and
\ref{morphisms-lemma-quasi-finite}).
By Morphisms, Lemma \ref{morphisms-lemma-generically-finite}
for every $\eta \in U$ which is the generic point of an
irreducible component of $U$, there exists an open neighbourhood
$V \subset U$ of $\eta$ such that $s^{-1}(V) \to V$ is finite. By
Descent, Lemma \ref{descent-lemma-descending-property-finite}
being finite is fpqc (and in particular \'etale) local on the target.
Hence we may apply
More on Groupoids, Lemma \ref{more-groupoids-lemma-property-invariant}
which says that the largest open $W \subset U$ over which $s$ is
finite is $R$-invariant. By the above $W$ contains every
generic point of an irreducible component of $U$.
The restriction $R_W$ of $R$ to $W$ equals $R_W = s^{-1}(W) = t^{-1}(W)$
(see Groupoids, Definition \ref{groupoids-definition-invariant-open}
and discussion following it).
By construction $s_W, t_W : R_W \to W$ are finite \'etale.
Consider the open subspace $X' = W/R_W \subset X$ (see
Spaces, Lemma \ref{spaces-lemma-finding-opens}).
By construction the inclusion map $X' \to X$
induces a bijection on points of codimension $0$.
This reduces us to
Lemma \ref{lemma-quasi-separated-finite-etale-cover-dense-open-scheme}.
\end{proof}

\end{document}
